% !TeX root = ../papers/nc-p3.tex
\section{三交換子の判別三次式と巡回行列からの復元}
\label{sec:marked-three-commutators}
\begin{proposition}\label{prop:marked-mixed-discriminant}
単純印を持つ有限窓の $q=3$, $g=4$ の場合に、定理\ref{thm:marked-common-triple-symmetry}の
基底で $P=Q$ とし、跡零対称行列の五空間 $S$ 内の二つの三平面を $G,F$ とする。
$G+F=N^\perp$ と跡対合で書き、$G\cap F=\langle H_1,H_2\rangle$ とする。
交代添字 $s\wedge t$ に対応する三座標を $v=s\times t$ とすれば、固定した非零定数を除いて
\[
          \det[Q_s,Q_t]=
          \operatorname{tr}\!\left([v]_\times[H_1,H_2]\right)\, v^t N v .
\]
従って常に特異交換子の直線があり、恒等零となる必要十分条件は $H_1,H_2$ の可換性である。
重複因子、可約二次式、等方方向を許す。
\end{proposition}
\begin{proof}
$K=[v]_\times$ は外積行列とし、$S$ 上に $\omega_v(X,Y)=\operatorname{tr}(K[X,Y])$ を取る。
交換子式により $[Q_s,Q_t]$ の成分は $G,F$ の順序付き基底での混合行列の成分の一定倍である。
従って $G$ の基底を $H_1,H_2,X$、$F$ の基底を $H_1,H_2,Y$ に替えたとき
行列式は非零定数しか変わらない。混合行列式を展開すると
$\omega_v(H_1,H_2)$ と、$H_1,H_2,X,Y$ 上の $\omega_v$ のPfaffianの積になる。

$\omega_v$ の根基は $K$ と交換する跡零対称行列である。
$0\ne v$ なら $K$ は巡回行列である。実際 $v^tv\ne0$ なら固有値が相異なり、
$v^tv=0$ なら $K^2=vv^t\ne0$, $K^3=0$。従って根基は
$T_v=vv^t-(v^tv)I/3$ の直線である。
体積形式を固定して $\omega_v\wedge\omega_v/2$ を $S$ のベクトルとみなすと
各成分は $v$ の二次式で、根基に入る。$T_v$ の成分に共通因子はない
（非対角の三単項式だけで足りる）。よってそのベクトルは $T_v$ の非零定数倍である。
四平面でのPfaffianは従って $\operatorname{tr}(N T_v)=v^tNv$ の一定倍である。
$N\ne0$ と交代空間での跡対合の非退化性より最後の主張も従う。
\end{proof}

\begin{proposition}\label{prop:marked-three-commutators-reconstruction}
$q=3$, $g>3$ で三つの指定交換子が独立の場合、対称三列 $G,F$ を固定すれば、
適当な有限個の方向を使って、$Q$ の復元を跡零行列 $X$ 一つの二次以下の条件と
線型方程式の整合性で被覆できる。具体的に基底 $s,t,u$ で指定値を
$C=[Q_s,Q_t]$, $D=[Q_s,Q_u]$, $L=[Q_t,Q_u]$ と読む。
\begin{enumerate}
\item $X$ の最小多項式が三次で、$\operatorname{tr}XC=\operatorname{tr}XD
=\operatorname{tr}X^2 C=\operatorname{tr}X^2 D=0$ となる跡零行列を取る。
\item $[X,Y_0]=C$, $[X,Z_0]=D$ の跡零の一解ずつを取る。
$X'=X^2-\operatorname{tr}(X^2)I/3$ とすると、残りの式はただ一つの線型系
\[
 L-[Y_0,Z_0]= aD+b(XD+DX)-cC-d(XC+CX)
\]
である。全ての解から $Q_s=X$, $Q_t=Y_0+aX+bX'$, $Q_u=Z_0+cX+dX'$ を得る。
\end{enumerate}
これは指定した印付き交換子の全解のうち $Q$ が既約なものを覆う。
得られた行列には元の点窓および正則性の条件が別に必要である。
\end{proposition}
\begin{proof}
巡回 $X$ の中心化環は $I,X,X^2$ の張る可換環で、$[X,-]$ の像はその跡直交補である。
$C,D$ の跡は交換子なので既に零である（候補 $G,F$ から始めるなら最初に三つの指定値の
跡を検査する）。これで一解の存在と一般解の表示が従う。
両中心化元の交換子は零、$[X^2,Z_0]=XD+DX$ などを代入すれば表示の線型系となり、
これを解けば三つの指定値を全て復元する。

最後に方向を有限個で選べることを示す。三平面の全元 $X$ が非巡回なら、非零跡零元は
二つの固有値のうち重複固有値の固有空間が二次元以上か、平方零で階数1以下である。
$\operatorname{tr}(X^2)$ が恒等零なら全元が後者となる。それ以外では三平面上の多項式恒等式
$(\operatorname{tr}X^2)^3=6(\operatorname{tr}X^3)^2$ が成立する。
一意分解により非零定数を除いて二次式は一次式の平方、三次式はその立方であり、
重複固有値 $-\operatorname{tr}(X^3)/\operatorname{tr}(X^2)$ は一次式 $\lambda(X)$ に延びる。
従って前者でも $X-\lambda(X)I$ は全三平面で階数1以下である。階数1以下の線型空間は
像直線を共有するか核平面を共有するので、元の三列の既約性に反する。
従って $I,X,X^2$ の独立を検出する三次小行列の一つは $s$ の非零な三次式である。
例えば $s=(1,i,j)$, $0\le i,j\le3$ の十六方向を取ればどれかで非零となる。
各方向は一度だけ基底に補えばよい。
\end{proof}
